\section{Data, features and evaluation protocol}\label{sec:data}

\subsection{Universe}
We study nineteen daily series from January 2005 (or inception) to October 2026
(\cref{tab:desc}): four Chinese equity indices (CSI~300, CSI~500, SSE~50, ChiNext); six international
equity indices (S\&P~500, Nasdaq~100, Hang Seng, Nikkei~225, DAX, FTSE~100); seven futures
(CSI~300 index futures IF, SHFE gold and copper, DCE soybean meal, WTI crude, COMEX gold, US 10-year
T-note); Bitcoin; and the CSI~300 ETF (510300), which is the tradable instrument for the strategy in
\cref{sec:economic}. Chinese series are obtained through AkShare, international series through Yahoo
Finance; every file is stored with a SHA-256 hash in a manifest. Futures are main-contract continuous
series and contain roll gaps; ETF prices are not dividend-adjusted, so its buy-and-hold returns are
understated by the dividend yield. For CSI~300 we also have 5-minute bars from May 2009 to July 2025,
from which we compute realised variance $\mathrm{RV}_t=\sum_ir_{t,i}^2$, bipower variation
\citep{barndorff2004power}, jumps and realised semivariances \citep{patton2015good}. They serve as
variance proxies for evaluation \citep{andersen2003modeling} and, in \cref{sec:rv}, as covariates.
The vendor file repeats every bar on 36 days and contains one bar with a zero price; we drop
duplicated time stamps and non-positive prices before computing any measure (an earlier version
of our code did not, which left realised variance unchanged but set bipower variation to zero on the
duplicated days). The CBOE VIX, used as a candidate gate variable in \cref{sec:kgate}, is taken from
Yahoo Finance.

\subsection{Features}
All features at date $t$ use prices up to the close of $t$; the target is $y_t=r_{t+1}$ in percent.
Let $v_t$ be the \citet{parkinson1980extreme} range variance
$v_t=(100\log(H_t/L_t))^2/(4\log 2)$ and $\hat\sigma_t=(\frac1{22}\sum_{j<22}v_{t-j})^{1/2}$. Then
$\mathrm{mom}_h=\sum_{j<h}r_{t-j}/(\hat\sigma_t\sqrt h)$ for $h\in\{20,120\}$ (trend);
$\mathrm{rev}_1=r_t/\hat\sigma_t$, $\mathrm{rev}_5=\sum_{j<5}r_{t-j}/(\hat\sigma_t\sqrt5)$ (reversal);
$\log\bar v^{(h)}_t$ for $h\in\{1,5,22\}$ (variance); and gate states
$(\log\bar v^{(22)}_t,\mathrm{mom}_{60},\log(P_t/\max_{j<250}P_{t-j}))$. Covariates are standardised
with training-window moments at each re-estimation.

\subsection{Protocol}
\begin{enumerate}[leftmargin=1.5em]
\item \textbf{Walk-forward.} Expanding window; statistical models are re-estimated at the start of
every quarter, the deep model every year. Out-of-sample forecasts start in January 2012 (later for
series with short histories, which require 750 training observations).
\item \textbf{Leakage control.} A model fitted for the block starting on date $d$ sees only rows with
date $<d$, whose targets are realised by the close of the first forecast date. The test suite
corrupts every price after a cut-off and verifies that all forecasts up to the cut-off are unchanged
for every model (\cref{app:leakage}).
\item \textbf{Lockbox.} Observations from 1 July 2024 onward were not used in any modelling decision.
All development results end on 30 June 2024; the lockbox period is evaluated once, after the protocol
was frozen, and reported separately.
\item \textbf{Common sample.} Within each asset all models are compared on the intersection of their
forecast dates.
\end{enumerate}

\subsection{Models}
\Cref{tab:models} lists the forecasters. All return a Gaussian mixture predictive distribution;
Student-$t$ GARCH densities are represented as Gaussian scale mixtures through 64-node generalised
Gauss--Laguerre quadrature (relative density error below $0.11\%$ for $\nu\ge3$; \cref{app:proofs}),
so every model is scored by the same closed-form functionals.

\begin{table}[t]\centering\small
\caption{Forecasting models}\label{tab:models}
\begin{adjustbox}{max width=\textwidth}\begin{tabular}{lll}\toprule
Model & Description & Refit\\\midrule
HistGauss & $\N$(250-day mean, 250-day variance) & quarterly\\
GARCH-t & constant mean, GARCH(1,1), Student-$t$ \citep{bollerslev1986generalized} & quarterly\\
AR-GJR-t & AR(1) mean, GJR-GARCH(1,1) \citep{glosten1993relation}, Student-$t$ & quarterly\\
Gauss-HAR & SMoE with $K=1$ (all mean covariates, log-HAR variance) & quarterly\\
MAR & SMoE with constant gate \citep{wong2001mixture} & quarterly\\
LSTAR & trend/reversal experts, logistic gate in volatility & quarterly\\
MS-AR & 2-regime Markov-switching regression \citep{hamilton1989new} & quarterly\\
XGB-Gauss & gradient-boosted mean, log-HAR residual variance & quarterly\\
\textbf{SMoE} & \cref{def:smoe}, $K=3$ & quarterly\\
Deep-SMoE & LSTM/MLP/CNN experts, attention gate, 3-seed ensemble & yearly\\
Chronos-Bolt & zero-shot foundation model \citep{ansari2024chronos}, small and base, Gaussian projection & none\\
TimesFM-2.5 & zero-shot foundation model \citep{das2024decoder}, 200M, Gaussian projection & none\\
\midrule
RealGARCH-t & log-linear Realized GARCH(1,1) \citep{hansen2012realized}, Student-$t$; \cref{sec:rv} only & quarterly\\
\bottomrule\end{tabular}\end{adjustbox}
\end{table}

\noindent The two foundation-model families differ in architecture and training. Chronos-Bolt is an
encoder--decoder T5 model released in November 2024 that outputs nine quantiles directly; TimesFM~2.5
is a decoder-only patch transformer released in September 2025 whose continuous quantile head we use
for the same nine levels. Both receive the last 512 daily returns as context and forecast one step
ahead without any fine-tuning. Their quantiles are projected onto a Gaussian by least squares on the
normal quantile function, so their tails beyond the 10\% and 90\% levels are extrapolated.

